\subsection{The exact reserve required by the weighted point sum}
\label{sec:archimedean-ap7-reserve}

The sufficient pointwise theorem
\ref{thm:archimedean-pointwise-sufficient} is enough for every
shift used in the positive-minor proof of (AP7).  In particular,
that proof does not require a pointwise theorem for every
near-parameter range allowed by the stronger source interface.
All norms, orbital integrals, character scales and measures in
the argument below are those in the original \(\mathrm{BData}\).

Write \(s=r+1\), \(m=2s\), and \(n=2r+k+1\).  Let
\(C_q\ge200\) be the fixed constant in the already established
moment theorem, and let \(C_g^{\rm mom}\) be its fixed all-gap
threshold.  The original definition is
\[
 q_0=\frac1{C_q}\min\left\{
        \frac{h_{NC}}{C_q\sqrt m},\ \operatorname{cloudGap}\right\}.
                                                        \tag{AR1}
\]
The cloud-gap set is retained as a nonempty finite set, as in
the actual weighted-sum consumer.  Its minimum is therefore
bounded below by any lower bound for all spectral gaps.

\begin{proposition}\label{prop:archimedean-ap7-reserve}
Fix the requested weight \(a_w\ge0\), and take the pointwise
constant supplied by Theorem~\ref{thm:archimedean-pointwise-sufficient},
for example
\[
                 C_*=64+16\max\{C_h,0\}+24C_J,
\]
where \(C_J\ge0\) is its fixed moment constant.  Set
\[
                  M=\lceil2a_w+4C_*+8\rceil,
                      P=Ms^3 .                         \tag{AR2}
\]
Choose the internal analytic base threshold \(\Gamma_g\) so that
\[
 \Gamma_g\ge\max\left\{C_{g,0},\ C_q^2(M+50),\
                   200(M+50),\ C_g^{\rm mom}(M+50)\right\},
                                                        \tag{AR3}
\]
where \(C_{g,0}\) denotes the other fixed local thresholds
in the sufficient pointwise theorem and the original construction.
Then the original strengthened base all-gap hypothesis
\(\operatorname{AllGaps}(\lambda_0,\Gamma_gm^4)\), together with
the original noncompact-gap hypothesis, supplies
\[
       q_0\ge s\bigl(P+48s^2+s+1\bigr).                 \tag{AR4}
\]
Consequently every shifted coefficient used in the positive-minor
proof of (AP7) satisfies the exact pointwise reserve.  The
weighted-sum estimate (AP7) follows without an additional
rank-dependent or shift-dependent hypothesis.
\end{proposition}

\begin{proof}
Put \(K=3sn^2+n^3\ge1\).  The original noncompact-gap
hypothesis and the strengthened all-gap hypothesis give
\[
 \frac{h_{NC}}{C_q\sqrt m}\ge\frac{\Gamma_gm^4K}{C_q},
       \qquad \operatorname{cloudGap}\ge \Gamma_gm^4.
\]
Using \(C_q\ge1\) in (AR1), we obtain
\[
 q_0\ge\frac{\Gamma_gm^4}{C_q^2}\min\{K,C_q\}
      \ge\frac{\Gamma_gm^4}{C_q^2}
      \ge16(M+50)s^4.                                  \tag{AR5}
\]
For \(s\ge1\),
\[
 s(P+48s^2+s+1)
      =Ms^4+48s^3+s^2+s\le(M+50)s^4,
\]
which proves (AR4).  The harmless factor \(16\) in
(AR5) is unused.

In the proof of (AP7), each selected positive-minor monomial
has total exponent at most \(P\).  Its exterior-Pieri shifts
therefore satisfy \(\delta=\|\lambda'-\lambda_0\|_1\le sP\).
The sufficient pointwise proof uses the Gram power \(p_0=6s^2\)
and a compact Wigner chain of length \(a\le36s^2\).  Thus
\[
             L=2p_0+a\le48s^2,
                  d=sL\le48s^3 .                      \tag{AR6}
\]
Its required moment reserve follows directly from (AR4):
\[
       \delta+s(L+s+1)
             \le s(P+48s^2+s+1)\le q_0.                \tag{AR7}
\]
This simultaneously contains the smaller reserve
\(s(P+s+1)\le q_0\) for the positive-minor iteration itself.

The all-gap assumptions are also simultaneous.  Since
\(Pm=2Ms^4\), \(Lm\le96s^3\), and
\(\delta+d\le(M+48)s^4\), (AR3) gives
\[
 \begin{split}
 \Gamma_gm^4&\ge C_g^{\rm mom}\max\{Pm,Lm\},\\
 \Gamma_gm^4&\ge200(\delta+d).
 \end{split}                                           \tag{AR8}
\]
In particular the required lower bound \(200\delta\)
holds, and all signed standard and dual-standard shifts in
the pointwise proof stay in the original regular chamber.
The original compact data and the base parameter are the
ones used for the spectral normalizer; they are unchanged
as the coefficient parameter varies.

Apply Theorem~\ref{thm:archimedean-pointwise-sufficient} to
every final coefficient in the positive-minor expansion.
It supplies exactly (AP6), with the same \(C_*\), for the
family actually used in that expansion.  The proof of
Theorem~\ref{thm:rotated-sum-from-near-bp} then gives
\[
 BSumR\bigl(\gamma\mapsto normWt(\gamma)^{a_ws^2}\bigr)
       \le b^{C_{\rm rot}m^5f}\,jSqG.
                                                        \tag{AR9}
\]
Its positive-minor selection, global product-formula
cancellation, and reciprocal-weight sum require no estimate
outside this family.  The constants in (AR3) and (AR9)
depend on the requested fixed weight and the original
fixed analytic constants, and not on \(r,k,f\).
\end{proof}

Together with the exact uniform amplitude and row-norm
comparisons, (AR9) supplies the weighted analytic input
used by \(\mathrm{StepThmBPinnedG}\).  This conclusion uses the fourth-power all-gap hypothesis in the
consumer's analytic reduction. Its internal fixed threshold
\(\Gamma_g\) is distinct from the externally prescribed
\(C_g\): the original existential representation-size and
level thresholds are chosen afterward to supply it.
The literal quantifier argument is given in
Theorem~\ref{thm:archimedean-pinned-row-closure}.
It does not assert the stronger full near-parameter
\(\mathrm{ThmBpDS}\) statement outside (AR7).
